\chapter{二次型}

\section{逐题解析}

\subsection{6.1 二次型的定义及其表示}

\begin{solutionblock}
\textbf{6.1 例 1}

\question{求二次型 \(f=x_1^2+3x_2^2-x_3^2+2x_1x_2+2x_1x_3+6x_2x_3\) 的矩阵及二次型的秩。}



\analysis{二次型
\[
f=x_1^2+3x_2^2-x_3^2+2x_1x_2+2x_1x_3+6x_2x_3
\]
写成 \(f=x^TAx\) 时，矩阵 \(A\) 必须为对称矩阵。平方项系数直接放在主对角线，交叉项 \(2a_{ij}x_ix_j\) 的系数等于题中交叉项系数。

因此
\[
a_{11}=1,\quad a_{22}=3,\quad a_{33}=-1,
\]
\[
2a_{12}=2,\quad 2a_{13}=2,\quad 2a_{23}=6.
\]
故
\[
A=\begin{pmatrix}
1&1&1\\
1&3&3\\
1&3&-1
\end{pmatrix}.
\]
矩阵的秩为
\[
r(A)=3.
\]}
\answer{\[
A=\begin{pmatrix}
1&1&1\\
1&3&3\\
1&3&-1
\end{pmatrix},\qquad r(f)=r(A)=3.
\]}
\end{solutionblock}

\begin{solutionblock}
\textbf{6.1 例 2}

\question{设 \(f(x_1,x_2,x_3)=\displaystyle\sum_{i=1}^3\sum_{j=1}^3ijx_ix_j\)，求该二次型对应的矩阵 \(A\)。}



\analysis{题中
\[
f(x_1,x_2,x_3)=\sum_{i=1}^3\sum_{j=1}^3 ijx_ix_j.
\]
所以二次型矩阵的第 \(i\) 行第 \(j\) 列元素为
\[
a_{ij}=ij.
\]
列出矩阵：
\[
A=\begin{pmatrix}
1\cdot1&1\cdot2&1\cdot3\\
2\cdot1&2\cdot2&2\cdot3\\
3\cdot1&3\cdot2&3\cdot3
\end{pmatrix}
=\begin{pmatrix}
1&2&3\\
2&4&6\\
3&6&9
\end{pmatrix}.
\]}
\answer{\[
A=\begin{pmatrix}
1&2&3\\
2&4&6\\
3&6&9
\end{pmatrix}.
\]}
\end{solutionblock}

\begin{solutionblock}
\textbf{6.1 例 3}

\question{设 \(f(x_1,x_2,x_3)=x^T\begin{pmatrix}1&4&0\\0&4&0\\0&0&0\end{pmatrix}x\)，求该二次型的秩。}



\analysis{给出的矩阵
\[
M=\begin{pmatrix}
1&4&0\\
0&4&0\\
0&0&0
\end{pmatrix}
\]
不是对称矩阵。二次型
\[
x^TMx
\]
只由矩阵的对称部分决定，即
\[
x^TMx=x^T\left(\frac{M+M^T}{2}\right)x.
\]
因此真正对应的二次型矩阵为
\[
A=\frac{M+M^T}{2}
=\begin{pmatrix}
1&2&0\\
2&4&0\\
0&0&0
\end{pmatrix}.
\]
前两行成比例，且第一行非零，故
\[
r(A)=1.
\]}
\answer{\(1\)}
\examnote{二次型矩阵默认取对称矩阵。若题面给了非对称矩阵，秩要算其对称部分的秩，而不是原矩阵的秩。}
\end{solutionblock}

\begin{solutionblock}
\textbf{6.1 例 4}

\question{设 \(f=(x_1+x_2+x_3)^2+(x_2-x_3)^2+(x_1+2x_2+x_3)^2\)，求二次型矩阵 \(A\)。}



\analysis{令
\[
l_1=x_1+x_2+x_3,\quad l_2=x_2-x_3,\quad l_3=x_1+2x_2+x_3.
\]
则
\[
f=l_1^2+l_2^2+l_3^2.
\]
把线性形式写成矩阵
\[
L=\begin{pmatrix}
1&1&1\\
0&1&-1\\
1&2&1
\end{pmatrix},
\qquad
\begin{pmatrix}l_1\\l_2\\l_3\end{pmatrix}=Lx.
\]
于是
\[
f=(Lx)^T(Lx)=x^T(L^TL)x,
\]
故二次型矩阵为
\[
A=L^TL
=\begin{pmatrix}
2&3&2\\
3&6&2\\
2&2&3
\end{pmatrix}.
\]}
\answer{\[
A=\begin{pmatrix}
2&3&2\\
3&6&2\\
2&2&3
\end{pmatrix}.
\]}
\end{solutionblock}

\subsection{6.2 可逆线性变换与矩阵合同}

\begin{solutionblock}
\textbf{6.2 例 1}

\question{设 \(\alpha_1=(1,2)^T,\alpha_2=(a,1)^T,x=(x_1,x_2)^T\)，令 \(f=\sum_{i=1}^2(\alpha_i,x)^2\)。若 \(f\) 与 \(g=y_1^2+y_2^2+2y_1y_2\) 合同，求 \(a\)。}



\analysis{由
\[
\alpha_1=(1,2)^T,\qquad \alpha_2=(a,1)^T
\]
得
\[
f=(\alpha_1^Tx)^2+(\alpha_2^Tx)^2=x^TAx,
\]
其中
\[
A=\alpha_1\alpha_1^T+\alpha_2\alpha_2^T
=\begin{pmatrix}
1+a^2&2+a\\
2+a&5
\end{pmatrix}.
\]
题中 \(f\) 经可逆线性变换化为
\[
g=y_1^2+y_2^2+2y_1y_2,
\]
其矩阵为
\[
B=\begin{pmatrix}1&1\\1&1\end{pmatrix},
\]
秩为 \(1\)。合同变换保持矩阵秩，因此
\[
r(A)=r(B)=1.
\]
对 \(A\) 计算行列式：
\[
|A|=5(1+a^2)-(2+a)^2=(2a-1)^2.
\]
令 \(|A|=0\)，得
\[
a=\frac12.
\]}
\answer{\(\dfrac12\)}
\end{solutionblock}

\begin{solutionblock}
\textbf{6.2 例 2}

\question{设 \(A=\begin{pmatrix}1&2\\2&1\end{pmatrix}\)，\(B=\begin{pmatrix}1&4\\1&1\end{pmatrix}\)。判断是否存在可逆矩阵 \(C\) 使 \(C^TAC=B\)，以及是否存在正交矩阵 \(Q\) 使 \(Q^TAQ=B\)。}



\analysis{给定
\[
A=\begin{pmatrix}1&2\\2&1\end{pmatrix},\qquad
B=\begin{pmatrix}1&4\\1&1\end{pmatrix}.
\]

第 (1) 个说法：若存在可逆矩阵 \(C\)，使 \(C^TAC=B\)，则 \(B\) 必须是对称矩阵，因为 \(A=A^T\) 时
\[
(C^TAC)^T=C^TA^TC=C^TAC.
\]
但题中 \(B\) 不是对称矩阵，所以第 (1) 个说法错误。

第 (2) 个说法：若存在正交矩阵 \(Q\)，使 \(Q^TAQ=B\)，则 \(B\) 不仅应为对称矩阵，而且与 \(A\) 正交相似，特征值也应相同。题中 \(B\) 非对称，故第 (2) 个说法也错误。}
\answer{两个说法均错误。}
\end{solutionblock}

\subsection{6.3 二次型的标准形和规范形}

\begin{solutionblock}
\textbf{6.3 例 1}

\question{用正交变换 \(x=Qy\) 化二次型 \(f=x_1^2+x_2^2+x_3^2+4x_1x_2+4x_1x_3+4x_2x_3\) 为标准形。}



\analysis{二次型
\[
f=x_1^2+x_2^2+x_3^2+4x_1x_2+4x_1x_3+4x_2x_3
\]
的矩阵为
\[
A=\begin{pmatrix}
1&2&2\\
2&1&2\\
2&2&1
\end{pmatrix}.
\]
该矩阵的特征值为 \(5,-1,-1\)。对应 \(\lambda=5\) 的单位特征向量可取
\[
q_1=\frac1{\sqrt3}(1,1,1)^T.
\]
对应 \(\lambda=-1\) 的特征子空间为 \(x_1+x_2+x_3=0\)，可取一组正交单位向量
\[
q_2=\frac1{\sqrt2}(1,-1,0)^T,\qquad
q_3=\frac1{\sqrt6}(1,1,-2)^T.
\]
令
\[
Q=(q_1,q_2,q_3),
\]
则 \(Q\) 为正交矩阵，且
\[
Q^TAQ=\operatorname{diag}(5,-1,-1).
\]}
\answer{取
\[
Q=\begin{pmatrix}
\frac1{\sqrt3}&\frac1{\sqrt2}&\frac1{\sqrt6}\\
\frac1{\sqrt3}&-\frac1{\sqrt2}&\frac1{\sqrt6}\\
\frac1{\sqrt3}&0&-\frac2{\sqrt6}
\end{pmatrix},
\]
则在正交变换 \(x=Qy\) 下，
\[
f=5y_1^2-y_2^2-y_3^2.
\]}
\end{solutionblock}

\begin{solutionblock}
\textbf{6.3 例 2}

\question{用配方法化二次型 \(f=2x_1^2+3x_2^2+x_3^2+4x_1x_2-4x_1x_3-8x_2x_3\) 为标准形，并写出变换 \(x=Cy\)。}



\analysis{对
\[
f=2x_1^2+3x_2^2+x_3^2+4x_1x_2-4x_1x_3-8x_2x_3
\]
配方。先含 \(x_1\) 的项：
\[
2x_1^2+4x_1x_2-4x_1x_3
=2(x_1+x_2-x_3)^2-2(x_2-x_3)^2.
\]
故
\[
f=2(x_1+x_2-x_3)^2+\bigl[3x_2^2+x_3^2-8x_2x_3-2(x_2-x_3)^2\bigr].
\]
括号内化简为
\[
x_2^2-4x_2x_3-x_3^2=(x_2-2x_3)^2-5x_3^2.
\]
令
\[
y_1=x_1+x_2-x_3,\qquad y_2=x_2-2x_3,\qquad y_3=x_3.
\]
则
\[
f=2y_1^2+y_2^2-5y_3^2.
\]
由上式反解
\[
x_3=y_3,\qquad x_2=y_2+2y_3,\qquad x_1=y_1-y_2-y_3.
\]
所以
\[
x=Cy,\qquad
C=\begin{pmatrix}
1&-1&-1\\
0&1&2\\
0&0&1
\end{pmatrix}.
\]}
\answer{\[
f=2y_1^2+y_2^2-5y_3^2,\qquad
C=\begin{pmatrix}
1&-1&-1\\
0&1&2\\
0&0&1
\end{pmatrix}.
\]}
\end{solutionblock}

\begin{solutionblock}
\textbf{6.3 例 3}

\question{用配方法化二次型 \(f=2x_1x_2+4x_1x_3\) 为标准形，并写出变换 \(x=Cy\)。}



\analysis{题中
\[
f=2x_1x_2+4x_1x_3=2x_1(x_2+2x_3).
\]
令
\[
y_1=x_1+x_2+2x_3,\qquad
y_2=x_1-x_2-2x_3,\qquad
y_3=x_3.
\]
则
\[
y_1^2-y_2^2
=(x_1+x_2+2x_3)^2-(x_1-x_2-2x_3)^2
=4x_1(x_2+2x_3).
\]
因此
\[
f=\frac12y_1^2-\frac12y_2^2.
\]
反解得
\[
x_1=\frac12y_1+\frac12y_2,\qquad
x_2=\frac12y_1-\frac12y_2-2y_3,\qquad
x_3=y_3.
\]
所以
\[
x=Cy,\qquad
C=\begin{pmatrix}
\frac12&\frac12&0\\
\frac12&-\frac12&-2\\
0&0&1
\end{pmatrix}.
\]}
\answer{\[
f=\frac12y_1^2-\frac12y_2^2,\qquad
C=\begin{pmatrix}
\frac12&\frac12&0\\
\frac12&-\frac12&-2\\
0&0&1
\end{pmatrix}.
\]}
\end{solutionblock}

\begin{solutionblock}
\textbf{6.3 例 4}

\question{求二次型 \(f=x_1x_2+x_1x_3+2x_2x_3\) 的标准形，并给出相应可逆线性变换。}



\analysis{二次型
\[
f=x_1x_2+x_1x_3+2x_2x_3
\]
的矩阵为
\[
A=\begin{pmatrix}
0&\frac12&\frac12\\
\frac12&0&1\\
\frac12&1&0
\end{pmatrix}.
\]
作可逆变换
\[
x=Cy,\qquad
C=\begin{pmatrix}
1&1&-\sqrt2\\
1&-1&-\frac1{\sqrt2}\\
0&0&\frac1{\sqrt2}
\end{pmatrix}.
\]
直接计算得
\[
C^TAC=\operatorname{diag}(1,-1,-1).
\]
因此规范形为
\[
z_1^2-z_2^2-z_3^2.
\]
若不缩放第三个变量，也可先得到标准形
\[
y_1^2-y_2^2-2y_3^2.
\]}
\answer{规范形为
\[
z_1^2-z_2^2-z_3^2,
\]
可取上述 \(C\) 作为坐标变换矩阵。}
\end{solutionblock}

\begin{solutionblock}
\textbf{6.3 例 5}

\question{求二次型 \(f=x_1^2+3x_2^2+x_3^2+2x_1x_2+2x_1x_3+2x_2x_3\) 的规范形。}



\analysis{二次型
\[
f=x_1^2+3x_2^2+x_3^2+2x_1x_2+2x_1x_3+2x_2x_3
\]
对应矩阵
\[
A=\begin{pmatrix}
1&1&1\\
1&3&1\\
1&1&1
\end{pmatrix}.
\]
计算其特征值为
\[
4,\ 1,\ 0.
\]
正惯性指数等于正特征值个数，为 \(2\)；负惯性指数等于负特征值个数，为 \(0\)。}
\answer{正惯性指数为 \(2\)，负惯性指数为 \(0\)。}
\end{solutionblock}

\begin{solutionblock}
\textbf{6.3 例 6}

\question{求二次型 \(f=-2x_1x_2-2x_1x_3+6x_2x_3\) 的正惯性指数，并在选项 \(3,2,1,0\) 中作出判断。}



\analysis{二次型
\[
f=-2x_1x_2-2x_1x_3+6x_2x_3
\]
对应矩阵为
\[
A=\begin{pmatrix}
0&-1&-1\\
-1&0&3\\
-1&3&0
\end{pmatrix}.
\]
其特征值为
\[
-3,\quad \frac{3-\sqrt{17}}2,\quad \frac{3+\sqrt{17}}2.
\]
其中
\[
-3<0,\qquad \frac{3-\sqrt{17}}2<0,\qquad \frac{3+\sqrt{17}}2>0.
\]
所以正惯性指数为 \(1\)。}
\answer{选 C，正惯性指数为 \(1\)。}
\end{solutionblock}

\begin{solutionblock}
\textbf{6.3 例 7}

\question{求二次型 \(f=(x_1+x_2)^2+(x_2-x_3)^2+(x_3+x_1)^2\) 的正惯性指数，并在选项 \(3,2,1,0\) 中作出判断。}



\analysis{因为
\[
f=(x_1+x_2)^2+(x_2-x_3)^2+(x_3+x_1)^2
\]
是三个平方和，所以负惯性指数为 \(0\)。展开可得矩阵
\[
A=\begin{pmatrix}
2&1&1\\
1&2&-1\\
1&-1&2
\end{pmatrix}.
\]
其特征值为 \(3,3,0\)，因此正惯性指数为 \(2\)。}
\answer{选 B，正惯性指数为 \(2\)。}
\end{solutionblock}

\subsection{6.4 正定二次型}

\begin{solutionblock}
\textbf{6.4 例 1}

\question{判断二次型 \(f=2x_1^2+4x_2^2+5x_3^2-4x_1x_3\) 是否正定。}



\analysis{二次型
\[
f=2x_1^2+4x_2^2+5x_3^2-4x_1x_3
\]
对应矩阵为
\[
A=\begin{pmatrix}
2&0&-2\\
0&4&0\\
-2&0&5
\end{pmatrix}.
\]
用顺序主子式判别：
\[
\Delta_1=2>0,
\]
\[
\Delta_2=\begin{vmatrix}2&0\\0&4\end{vmatrix}=8>0,
\]
\[
\Delta_3=|A|=24>0.
\]
三个顺序主子式全为正，故 \(A\) 正定，二次型 \(f\) 正定。}
\answer{\(f\) 是正定二次型。}
\end{solutionblock}

\begin{solutionblock}
\textbf{6.4 例 2}

\question{求使二次型 \(f=2x_1^2+x_2^2+x_3^2+2x_1x_2+tx_2x_3\) 正定的参数 \(t\) 的取值范围。}



\analysis{二次型
\[
f=2x_1^2+x_2^2+x_3^2+2x_1x_2+tx_2x_3
\]
对应矩阵为
\[
A=\begin{pmatrix}
2&1&0\\
1&1&\frac t2\\
0&\frac t2&1
\end{pmatrix}.
\]
顺序主子式为
\[
\Delta_1=2>0,
\]
\[
\Delta_2=\begin{vmatrix}2&1\\1&1\end{vmatrix}=1>0,
\]
\[
\Delta_3=|A|=1-\frac{t^2}{2}.
\]
正定要求 \(\Delta_3>0\)，所以
\[
1-\frac{t^2}{2}>0
\quad\Longleftrightarrow\quad
t^2<2.
\]}
\answer{\[
-\sqrt2<t<\sqrt2.
\]}
\end{solutionblock}

\begin{solutionblock}
\textbf{6.4 例 3}

\question{设 \(A=\begin{pmatrix}1&0&1\\0&2&0\\1&0&1\end{pmatrix}\)，\(B=(A+kE)^2\)。求使 \(B\) 正定的 \(k\) 的取值范围。}



\analysis{矩阵
\[
A=\begin{pmatrix}
1&0&1\\
0&2&0\\
1&0&1
\end{pmatrix}
\]
为实对称矩阵，特征值为
\[
0,\ 2,\ 2.
\]
矩阵
\[
B=(A+kE)^2
\]
的特征值为
\[
(0+k)^2,\quad (2+k)^2,\quad (2+k)^2.
\]
\(B\) 正定的充要条件是这三个特征值均大于 \(0\)，即
\[
k\ne0,\qquad k\ne-2.
\]}
\answer{\[
k\in\mathbb{R}\setminus\{0,-2\}.
\]}
\end{solutionblock}

\begin{solutionblock}
\textbf{6.4 例 4}

\question{设 \(f=(ax_1+2x_2-3x_3)^2+(x_2-2x_3)^2+(x_1+ax_2-x_3)^2\)，判断 \(f\) 正定时 \(a\) 应满足的条件。}



\analysis{题中
\[
f=(ax_1+2x_2-3x_3)^2+(x_2-2x_3)^2+(x_1+ax_2-x_3)^2.
\]
这是三个线性形式平方和。设
\[
R=\begin{pmatrix}
a&2&-3\\
0&1&-2\\
1&a&-1
\end{pmatrix}.
\]
则
\[
f=\|Rx\|^2=x^TR^TRx.
\]
该二次型正定的充要条件是 \(Rx=0\) 只有零解，即 \(R\) 可逆。计算
\[
|R|=(a-1)(2a+1).
\]
故正定要求
\[
(a-1)(2a+1)\ne0,
\]
即
\[
a\ne1,\qquad a\ne-\frac12.
\]}
\answer{选 C，\(a\ne1\) 且 \(a\ne-\dfrac12\)。}
\end{solutionblock}

\begin{solutionblock}
\textbf{6.4 例 5}

\question{设 \(f=x^TAx\)，\(A^T=A\)。判断 \(f\) 正定的充要条件。}



\analysis{对称矩阵 \(A\) 对应二次型 \(f=x^TAx\)。若 \(f\) 正定，则 \(A\) 正定，必可逆，因此 \(A^{-1}\) 也正定；反之若 \(A^{-1}\) 正定，则其逆矩阵 \(A\) 也正定，故 \(f\) 正定。因此 D 是充要条件。

选项 A 只有 \(|A|>0\)，不足以保证所有特征值均为正。选项 B “负惯性指数为 \(0\)”只说明没有负平方项，但可能存在零平方项，属于半正定。选项 C “秩为 \(n\)”只说明非退化，也可能是不定型。}
\answer{选 D。}
\end{solutionblock}

\begin{solutionblock}
\textbf{6.4 例 6}

\question{设 \(A\) 为 \(m\times n\) 矩阵，\(E\) 为 \(n\) 阶单位矩阵，\(B=\lambda E+A^TA\)。证明当 \(\lambda>0\) 时 \(B\) 正定。}



\analysis{设 \(x\ne0\)。由
\[
B=\lambda E+A^TA
\]
得
\[
x^TBx=\lambda x^Tx+x^TA^TAx
=\lambda\|x\|^2+\|Ax\|^2.
\]
当 \(\lambda>0\) 时，
\[
\lambda\|x\|^2>0,\qquad \|Ax\|^2\ge0,
\]
所以
\[
x^TBx>0.
\]
故 \(B\) 为正定矩阵。}
\answer{已证：当 \(\lambda>0\) 时，\(B\) 为正定矩阵。}
\end{solutionblock}

\begin{solutionblock}
\textbf{6.4 例 7}

\question{设 \(A\) 为 \(n\) 阶正定矩阵，证明存在 \(n\) 阶可逆矩阵 \(D\) 使 \(A=D^TD\)，并证明存在 \(n\) 阶正定矩阵 \(B\) 使 \(A=B^2\)。}



\analysis{第 (1) 问：若 \(A\) 正定，则存在正交矩阵 \(Q\)，使
\[
Q^TAQ=\Lambda=\operatorname{diag}(\lambda_1,\ldots,\lambda_n),
\]
且
\[
\lambda_i>0.
\]
令
\[
\Lambda^{1/2}=\operatorname{diag}(\sqrt{\lambda_1},\ldots,\sqrt{\lambda_n}),
\qquad
D=\Lambda^{1/2}Q^T.
\]
则 \(D\) 可逆，且
\[
D^TD=Q\Lambda^{1/2}\Lambda^{1/2}Q^T=Q\Lambda Q^T=A.
\]

第 (2) 问：令
\[
B=Q\Lambda^{1/2}Q^T.
\]
则 \(B\) 为实对称矩阵，且特征值均为 \(\sqrt{\lambda_i}>0\)，所以 \(B\) 正定。并且
\[
B^2=Q\Lambda^{1/2}Q^TQ\Lambda^{1/2}Q^T
=Q\Lambda Q^T=A.
\]}
\answer{两个结论均成立。}
\examnote{这是正定矩阵的平方根分解：\(A=D^TD\) 是 Cholesky 型分解，\(A=B^2\) 是对称正定平方根。}
\end{solutionblock}

\subsection{6.5 矩阵等价、相似、合同}

\begin{solutionblock}
\textbf{6.5 例 1}

\question{设 \(A=\begin{pmatrix}2&-1&-1\\-1&2&-1\\-1&-1&2\end{pmatrix}\)，\(B=\begin{pmatrix}1&0&0\\0&1&0\\0&0&0\end{pmatrix}\)，判断 \(A\) 与 \(B\) 的等价、相似、合同关系。}



\analysis{设
\[
A=\begin{pmatrix}
2&-1&-1\\
-1&2&-1\\
-1&-1&2
\end{pmatrix},\qquad
B=\begin{pmatrix}
1&0&0\\
0&1&0\\
0&0&0
\end{pmatrix}.
\]
矩阵 \(A\) 的特征值为
\[
3,\ 3,\ 0,
\]
因此 \(r(A)=2\)。而 \(r(B)=2\)，所以 \(A\) 与 \(B\) 等价。

相似要求特征值完全相同。\(B\) 的特征值为
\[
1,\ 1,\ 0,
\]
与 \(A\) 的特征值不同，所以二者不相似。

合同分类看惯性指数。\(A\) 的正惯性指数为 \(2\)、负惯性指数为 \(0\)；\(B\) 的正惯性指数为 \(2\)、负惯性指数为 \(0\)。因此 \(A\) 与 \(B\) 合同。}
\answer{选 B，合同但不相似。}
\end{solutionblock}

\begin{solutionblock}
\textbf{6.5 例 2}

\question{设 \(A\) 为 \(n\) 阶实对称矩阵，将 \(A\) 的第 \(i\) 行与第 \(j\) 行交换，再将第 \(i\) 列与第 \(j\) 列交换得到 \(B\)，判断 \(A,B\) 的等价、相似、合同关系。}



\analysis{交换 \(A\) 的第 \(i\) 列和第 \(j\) 列，可写成右乘一个置换矩阵 \(P\)；再交换第 \(i\) 行和第 \(j\) 行，可写成左乘同一个置换矩阵 \(P\)。因此
\[
B=PAP.
\]
置换矩阵满足
\[
P^T=P,\qquad P^{-1}=P.
\]
于是
\[
B=P^TAP,
\]
所以 \(A\) 与 \(B\) 合同；同时
\[
B=P^{-1}AP,
\]
所以 \(A\) 与 \(B\) 相似。由于初等行列变换不改变等价关系，且相似或合同都蕴含等价矩阵同秩，因此 \(A\) 与 \(B\) 也等价。}
\answer{选 D，等价、相似、合同。}
\end{solutionblock}
